% INTERIM REPORT: the three sections that are reviewed first (Introduction, Literature Review,
% Problem Statement). The section files are shared with the full paper (main.tex), so editing
% a section updates both documents. Build: `uv run qdrift paper --build`.
\documentclass[10pt,conference]{IEEEtran}
% Shared by main.tex (full paper) and interim.tex (Introduction, Literature Review and Problem
% Statement only). Numbers come from generated/numbers.tex through \val{key}; tables and figures
% come from generated/. Never type a result by hand.
\IEEEoverridecommandlockouts
\usepackage[T1]{fontenc}
\usepackage{cite}
\usepackage{amsmath,amssymb}
\usepackage{array}
\usepackage{graphicx}
\usepackage{booktabs}
\usepackage{xcolor}
\usepackage[hidelinks]{hyperref}
\graphicspath{{generated/figures/}}

% Generated by `uv run qdrift paper`; do not edit.
\providecommand{\qdset}[2]{\expandafter\def\csname qd@#1\endcsname{#2}}
\qdset{acc-EMU-breastcancer-M1}{0.963}
\qdset{acc-EMU-breastcancer-M2}{0.907}
\qdset{acc-EMU-breastcancer-M3}{0.883}
\qdset{acc-EMU-twocurves-M1}{0.917}
\qdset{acc-EMU-twocurves-M2}{0.760}
\qdset{acc-EMU-twocurves-M3}{0.797}
\qdset{acc-R0-breastcancer-M1}{0.887}
\qdset{acc-R0-breastcancer-M2}{0.893}
\qdset{acc-R0-breastcancer-M3}{0.837}
\qdset{acc-R0-twocurves-M1}{0.813}
\qdset{acc-R0-twocurves-M2}{0.707}
\qdset{acc-R0-twocurves-M3}{0.647}
\qdset{acc-R1-breastcancer-M1}{0.893}
\qdset{acc-R1-breastcancer-M2}{0.893}
\qdset{acc-R1-breastcancer-M3}{0.807}
\qdset{acc-R1-twocurves-M1}{0.800}
\qdset{acc-R1-twocurves-M2}{0.663}
\qdset{acc-R1-twocurves-M3}{0.670}
\qdset{acc-R216-breastcancer-M1}{0.887}
\qdset{acc-R216-breastcancer-M2}{0.897}
\qdset{acc-R216-breastcancer-M3}{0.830}
\qdset{acc-R216-twocurves-M1}{0.803}
\qdset{acc-R216-twocurves-M2}{0.727}
\qdset{acc-R216-twocurves-M3}{0.663}
\qdset{acc-R216s1-breastcancer-M1}{0.887}
\qdset{acc-R216s1-breastcancer-M2}{0.897}
\qdset{acc-R216s1-breastcancer-M3}{0.830}
\qdset{acc-R216s1-twocurves-M1}{0.803}
\qdset{acc-R216s1-twocurves-M2}{0.727}
\qdset{acc-R216s1-twocurves-M3}{0.663}
\qdset{acc-R232-breastcancer-M1}{0.887}
\qdset{acc-R232-breastcancer-M2}{0.893}
\qdset{acc-R232-breastcancer-M3}{0.813}
\qdset{acc-R232-twocurves-M1}{0.810}
\qdset{acc-R232-twocurves-M2}{0.713}
\qdset{acc-R232-twocurves-M3}{0.643}
\qdset{acc-R28-breastcancer-M1}{0.883}
\qdset{acc-R28-breastcancer-M2}{0.893}
\qdset{acc-R28-breastcancer-M3}{0.830}
\qdset{acc-R28-twocurves-M1}{0.800}
\qdset{acc-R28-twocurves-M2}{0.703}
\qdset{acc-R28-twocurves-M3}{0.663}
\qdset{acc-R3m-breastcancer-M1}{0.897}
\qdset{acc-R3m-breastcancer-M2}{0.897}
\qdset{acc-R3m-breastcancer-M3}{0.860}
\qdset{acc-R3m-twocurves-M1}{0.827}
\qdset{acc-R3m-twocurves-M2}{0.727}
\qdset{acc-R3m-twocurves-M3}{0.667}
\qdset{acc-R3r-breastcancer-M1}{0.907}
\qdset{acc-R3r-breastcancer-M2}{0.907}
\qdset{acc-R3r-breastcancer-M3}{0.843}
\qdset{acc-R3r-twocurves-M1}{0.810}
\qdset{acc-R3r-twocurves-M2}{0.673}
\qdset{acc-R3r-twocurves-M3}{0.663}
\qdset{acc-R4-breastcancer-M1}{0.890}
\qdset{acc-R4-breastcancer-M2}{0.897}
\qdset{acc-R4-breastcancer-M3}{0.827}
\qdset{acc-R4-twocurves-M1}{0.817}
\qdset{acc-R4-twocurves-M2}{0.750}
\qdset{acc-R4-twocurves-M3}{0.683}
\qdset{acc-SIM-breastcancer-M1}{0.967}
\qdset{acc-SIM-breastcancer-M2}{0.933}
\qdset{acc-SIM-breastcancer-M3}{0.913}
\qdset{acc-SIM-twocurves-M1}{0.927}
\qdset{acc-SIM-twocurves-M2}{0.873}
\qdset{acc-SIM-twocurves-M3}{0.847}
\qdset{aci-gamma}{0.03}
\qdset{alpha}{0.1}
\qdset{base-acc-breastcancer-logreg}{0.940}
\qdset{base-acc-breastcancer-mlp}{0.960}
\qdset{base-acc-breastcancer-rbfsvm}{0.940}
\qdset{base-acc-twocurves-logreg}{0.733}
\qdset{base-acc-twocurves-mlp}{0.973}
\qdset{base-acc-twocurves-rbfsvm}{0.987}
\qdset{bootstrap}{2,000}
\qdset{budget-cap}{1200}
\qdset{campaign-cap}{400}
\qdset{campaign-first}{2026-10-27}
\qdset{campaign-last}{2026-10-27}
\qdset{coverage-target}{90\%}
\qdset{devices}{\texttt{ibm\_fez} and \texttt{ibm\_kingston}}
\qdset{emulator-shots}{2000}
\qdset{fallback}{\texttt{ibm\_marrakesh}}
\qdset{gap-max-pp}{22.7}
\qdset{gap-mean-pp}{11.3}
\qdset{gap-min-pp}{4.0}
\qdset{h1-decision}{supported}
\qdset{h1-hi}{0.105}
\qdset{h1-lo}{0.044}
\qdset{h1-mean}{0.074}
\qdset{h3-decision}{not supported}
\qdset{h3-hi}{0.029}
\qdset{h3-lo}{-0.001}
\qdset{h3-mean}{0.013}
\qdset{h4-decision}{not supported}
\qdset{m1-features}{4}
\qdset{m1-layers}{3}
\qdset{m1-params}{28}
\qdset{m1-restarts}{3}
\qdset{m1-steps}{200}
\qdset{m2-features}{12}
\qdset{m2-layers}{1}
\qdset{m3-features}{10}
\qdset{m3-layers}{2}
\qdset{n-cal}{32}
\qdset{n-hw-samples}{182}
\qdset{n-qubits}{4}
\qdset{n-test}{150}
\qdset{n-train}{240}
\qdset{n-val}{100}
\qdset{noise-aware-shots}{2,000}
\qdset{patch-buffer}{1}
\qdset{patches}{14}
\qdset{per-execution-us}{265}
\qdset{pilot-cap}{90}
\qdset{qpu-campaign}{41}
\qdset{qpu-pilot}{22}
\qdset{qpu-used}{63}
\qdset{queue-max-h}{1.9}
\qdset{queue-median-h}{1.0}
\qdset{readout-shots}{2000}
\qdset{refit-lambda}{0.25}
\qdset{seed}{20261026}
\qdset{session-usage-mean}{20.5}
\qdset{sessions-done}{2}
\qdset{sessions-per-device}{5}
\qdset{sessions-planned}{10}
\qdset{shots}{500}
\qdset{sim-acc-breastcancer-M1}{0.967}
\qdset{sim-acc-breastcancer-M2}{0.933}
\qdset{sim-acc-breastcancer-M3}{0.913}
\qdset{sim-acc-twocurves-M1}{0.927}
\qdset{sim-acc-twocurves-M2}{0.873}
\qdset{sim-acc-twocurves-M3}{0.847}

\DeclareRobustCommand{\val}[1]{\ifcsname qd@#1\endcsname\csname qd@#1\endcsname\else\textbf{??#1}\fi}
\newcommand{\todo}[1]{\textcolor{red}{[TODO: #1]}}

% \iffull{full-paper text}{interim-report text}: the first argument is used when the Methodology
% section exists (full paper), the second when only the first three sections are built.
\newcommand{\iffull}[2]{\ifcsname r@sec:methods\endcsname #1\else #2\fi}

% Author block shared by both documents. The affiliation is inferred from the e-mail domain
% (nst.rishihood.edu.in); confirm it.
\newcommand{\qdauthor}{%
\IEEEauthorblockN{Pankaj Baid}
\IEEEauthorblockA{Newton School of Technology, Rishihood University\\
Sonipat, Haryana, India\\
pankaj.baid2024@nst.rishihood.edu.in}}


\begin{document}

\title{Noise Drift of Simulator-Trained Quantum Classifiers on IBM Heron Processors:
Introduction, Literature Review and Problem Statement}

\author{\qdauthor}

\maketitle

\noindent\textit{Interim report. It contains the Introduction, the Literature Review and the
Problem Statement of the full paper. The abstract, methodology, results and discussion, and
conclusion are in the full paper.}

% ===================================================================== INTRODUCTION
% Target: about one column. Write in your own words; cite with \cite{key} (keys in
% literature/refs.bib). Cover, in this order:
%  1. QML classifiers are usually trained in noiseless simulation and then run on NISQ
%     hardware, whose noise differs from the simulator and changes from day to day.
%     \cite{wang2021nibp,heyraud2022noisykernels}
%  2. What is known: the one-off simulator-to-hardware gap and fixes for it
%     \cite{wang2022quantumnat,nation2021m3}; noise that varies over time breaks conformal
%     guarantees, shown on one day of data \cite{park2023qcp,spencer2026aqcp}; multi-day drift
%     and emulator trust were studied for QAOA, not classifiers \cite{qaoa2026drift}; small
%     hardware test sets mislead \cite{qbiobench2026}.
%  3. The gap in one or two sentences: no multi-day, multi-device measurement for classifiers
%     that also prices each repair in QPU seconds under a free-tier budget.
%  4. Research questions RQ1-RQ4 (copy their meaning from preregistration.md, not the wording).
%  5. Contributions as a short itemize list, each tied to a result section. Only claim what the
%     final results support.
%  6. One sentence that no quantum advantage is claimed; tuned classical baselines are reported
%     \cite{bowles2024benchmarking,schuld2022advantage}.
\section{Introduction}
\label{sec:intro}
\todo{Write the introduction.}

% Drafted with an AI assistant on 12 Nov 2026. Every statement about another paper was checked
% against its abstract or text (files in literature/pdfs/). Revise it in your own words, and
% keep that discipline when you edit: do not add a claim you have not read.
\section{Literature Review}
\label{sec:litreview}

This section reviews the work on which the study builds. It describes the quantum classifiers
that are run on hardware, how noise affects quantum models, the methods proposed to adapt to
or correct noise, ways to quantify uncertainty, and critical benchmarks. It ends with a
comparison of the closest studies and the gap that remains.

\subsection{Quantum classifiers on near-term hardware}
Early hardware demonstrations of quantum classification used a superconducting processor to
implement a variational classifier and a kernel estimator based on a quantum feature map
\cite{havlicek2019supervised}. Later work clarified what such models can represent. A model
that re-uploads its input several times can act as a universal classifier even on a single
qubit \cite{perezsalinas2020reuploading}. More generally, a circuit with data-encoding gates is
a partial Fourier series in the input, whose accessible frequencies are set by the encoding
\cite{schuld2021encoding}, and the standard circuit-based models, including re-uploading
circuits, can be written as linear models in a quantum feature space
\cite{jerbi2023beyondkernels}. Projecting the quantum state back to a classical representation
gives ``projected'' models, and the same work shows that, with training data available,
classical learners can stay competitive even on problems tailored to quantum computers
\cite{huang2021power}. A further family uses a fixed quantum circuit as a feature generator and
trains only a classical readout, as in quantum reservoir computing, which has been
demonstrated on superconducting devices \cite{yasuda2023qrc}. Our three classifiers belong to
these families: a re-uploading circuit (M1), projected single-qubit features (M2) and a fixed
random circuit with a linear readout (M3).

\subsection{Noise and the failure of quantum models}
Noise harms quantum models in several ways. For variational circuits, local Pauli noise makes
the gradient vanish exponentially with circuit depth, a noise-induced barren plateau
\cite{wang2021nibp}. Kernel methods face a related problem: kernel values can concentrate
exponentially around a constant, and noise is one of four identified sources, next to the
expressivity of the embedding, global measurements and entanglement
\cite{thanasilp2024concentration}. Decoherence also limits the number of independent features
that a noisy quantum kernel can extract, and acts as an implicit regulariser
\cite{heyraud2022noisykernels}. Remedies studied for kernels include tuning the bandwidth
\cite{shaydulin2022bandwidth}, although tuned quantum kernels can then resemble classical
radial-basis kernels \cite{bandwidth2025similarity}; training the embedding while mitigating
the effect of device noise \cite{hubregtsen2022embedding}; covariant kernels with a bit-flip
tolerance \cite{covariant2025concentration}; local and multi-scale kernels
\cite{localkernels2026}; and allocating a limited shot budget to the kernel entries that matter
most \cite{aqka2026}. These works concern training, kernel estimation and scaling with the
number of qubits. Our setting is different: the models are small, they are trained only in
simulation, and the question is how their predictions change when they are executed on hardware.

\subsection{Adapting models to device noise}
Several methods make a model robust to the noise of a specific device. QuantumNAT reported a gap
of more than 60\% between noise-free simulation and hardware for a four-class digit task on
IBMQ-Yorktown, found that noise acts on the measured outputs roughly as a scaling and a shift,
and countered it with post-measurement normalisation, noise injection during training and
quantisation of the outputs, improving accuracy by up to 43\% on six devices
\cite{wang2022quantumnat}. QuantumNAS searches jointly for a noise-robust circuit and qubit
mapping, and its authors note that a searched circuit must be evaluated on the machine right
after the search, because otherwise the noise characteristics drift \cite{wang2022quantumnas}.
QOC trains the parameters on the device with the parameter-shift rule and prunes the gradients
that noise has made unreliable \cite{wang2022qoc}. Across devices, a recent study transfers a
learned noise model from one IBM processor to another and finds that two-qubit gate error is
the main source of mismatch \cite{farib2026crossdevice}. Because superconducting processors
fluctuate, calibration has to be repeated, and real-time calibration frameworks aim to do so
without taking the processor offline \cite{realtimecal2023}.

\subsection{Error mitigation}
Error mitigation corrects expectation values after the fact, usually at the price of extra
circuits or shots. Readout errors, which bias every measured expectation value, can be mitigated
in a reduced subspace of the observed bit strings \cite{nation2021m3} or without a noise model,
by turning the bias into a multiplicative factor that is measured and removed
\cite{vandenberg2022trex}. For gate errors, zero-noise extrapolation and probabilistic error
cancellation are standard techniques for short circuits \cite{temme2017zne}; Clifford data
regression learns the correction from near-Clifford circuits that can be simulated classically
\cite{czarnik2021cdr}; and machine-learned mitigation lowers the cost of mitigation, with its
authors listing the drift of noise in hardware as an open question \cite{liao2024mlqem}. On a
free budget the extra runtime matters, because mitigation circuits compete with the experiment
for the same processor time. For this reason we test only repairs that need no additional
circuits, a small number of calibration circuits, or a few labelled hardware samples.

\subsection{Uncertainty and calibration of quantum predictions}
A classifier should also report how far its predictions can be trusted. Calibration measures such
as the expected calibration error were popularised for neural networks \cite{guo2017calibration},
and a recent library provides shot-based and ensemble-based uncertainty estimates for quantum
models and reports that variational classifiers tend to be overconfident \cite{quantumuq2026}.
Finite-shot evaluation by itself randomises a deployed model, and an audit of frozen quantum
kernel classifiers found that this moves thresholded metrics by about 0.02
\cite{fernandezbarrios2026conditional}. Conformal prediction turns a few shots of a quantum
model into prediction sets with a coverage guarantee \cite{park2023qcp}. The guarantee assumes
exchangeable data, which time-varying hardware noise can violate. Adaptive conformal inference
\cite{gibbs2021aci} lifts the assumption by updating the working miscoverage level online, and
AQCP applies it to quantum models, with hardware data from a single day and a multi-day
assessment left to future work \cite{spencer2026aqcp}.

\subsection{Benchmarking and the evidence for quantum advantage}
Claims that quantum models outperform classical ones need careful baselines. A large study of
12 popular models on 6 binary classification tasks, with 160 datasets, found that
out-of-the-box classical models outperform the quantum classifiers and that removing
entanglement often leaves performance unchanged or better \cite{bowles2024benchmarking}. The
same study counts that about 40\% of a sample of 55 arXiv papers claim that a quantum model
beats a classical one, whereas only 3 report the opposite. A related argument holds that a focus
on advantage restricts attention to a biased subset of models and datasets
\cite{schuld2022advantage}. Quantum convolutional networks look strong on phase-recognition
benchmarks, yet on such locally easy datasets they can be simulated classically
\cite{bermejo2026qcnn}, and an attribution audit of intrusion-detection benchmarks traced
apparent gains to classical preprocessing and regularisation \cite{gillani2026nids}. On
hardware, QBioBench evaluated kernel and variational classifiers on an IBM Nighthawk processor:
no quantum model reached the performance of its matched support vector machine, readout
mitigation changed the ROC-AUC by at most 0.015, and random subsets of the test pool showed
that the small test sets favoured by circuit budgets give unreliable estimates
\cite{qbiobench2026}. We take two lessons from this literature: report tuned classical
baselines, and use a test set that is fixed across sessions.

\subsection{Hardware studies of calibration, drift and parallel execution}
Few studies follow the same workload on hardware over days. The closest is a study of QAOA for
weighted MaxCut on three IBM Heron r2 processors, which combined a 54-pair hardware grid,
repeats of the same circuits 9.3 to 12.1 days later and a refresh experiment. Selecting qubits
with the full calibration improved the median approximation ratio by 0.74 percentage points; an
emulator built from the calibration reproduced the direction of this advantage in 35 of 54
comparisons but overstated its size, and the compiler effect kept its sign in 16 of 18 repeats
\cite{qaoa2026drift}. The authors conclude that calibration data are useful for finding
candidates and deciding when to refresh, but that scalar calibration scores do not predict
hardware performance well. Running many small circuits side by side raises throughput but can
add crosstalk, which has to be characterised and scheduled around \cite{murali2020crosstalk}.
The palloq method chooses which circuits to combine and where to place them to limit this
interference \cite{ohkura2022simultaneous}, and a study of parallel QML circuits on an IBM
processor measures how the occupancy of the processor affects the loss of the learning
algorithm \cite{parallelqml2024}. Our packing of \val{patches} circuit copies per job follows
this line of work, with a pilot run to check for crosstalk. All circuits are built and
transpiled with Qiskit \cite{javadiabhari2024qiskit}.

\begin{table*}[t]
\centering
\caption{Closest prior studies and what remains open. Statements about other studies follow
their abstracts or text.}
\label{tab:related}
\footnotesize
\setlength{\tabcolsep}{4pt}
\renewcommand{\arraystretch}{1.15}
\begin{tabular}{>{\raggedright\arraybackslash}p{2.5cm}>{\raggedright\arraybackslash}p{3.3cm}>{\raggedright\arraybackslash}p{3.3cm}>{\raggedright\arraybackslash}p{3.2cm}>{\raggedright\arraybackslash}p{3.9cm}}
\toprule
Study & Workload & Hardware and period & Correction or uncertainty method & What remains open \\
\midrule
QuantumNAT \cite{wang2022quantumnat} & Quantum neural networks, 8 classification tasks & 6 IBM devices & Normalisation, noise injection, quantisation & Variation across days is not reported \\
QuantumNAS \cite{wang2022quantumnas} & Joint search of circuit and qubit mapping & IBM devices & Search under a noise model & Noise drifts, so a searched circuit must be run right after the search \\
AQCP \cite{spencer2026aqcp} & Probabilistic regression, 5 qubits & \texttt{ibm\_sherbrooke}, one day & Adaptive conformal inference & Multi-day assessment left to future work \\
QAOA drift study \cite{qaoa2026drift} & Weighted MaxCut with QAOA & 3 Heron r2 processors, repeats after 9--12 days & Calibration-aware qubit selection and refresh & Optimisation, not classification; emulator overstated the gain \\
QBioBench \cite{qbiobench2026} & Biosignal classification, kernels and variational circuits & One IBM Nighthawk r2 processor & Readout mitigation & Single processor, exploratory; small test sets unreliable \\
Cross-device transfer \cite{farib2026crossdevice} & Learning noise models of circuit outputs & \texttt{ibm\_fez} and \texttt{ibm\_marrakesh} & Few-shot fine-tuning, 20 samples & Targets output distributions, not classifier accuracy \\
Learned mitigation \cite{liao2024mlqem} & Expectation values, up to 100 qubits & IBM devices & Machine-learned mitigation & Drift of noise named as an open question \\
\midrule
\textbf{This study} & Three classifiers, two datasets & Two Heron r2 processors, \val{sessions-done} sessions over \val{campaign-days} days & Four repair families and four conformal schemes, each priced in QPU seconds & Frozen small models; free-plan budget \\
\bottomrule
\end{tabular}
\end{table*}

\subsection{Summary and research gap}
\label{sec:gap}
Table~\ref{tab:related} compares the closest studies. Four points emerge. First, evidence on how
deployed classifiers change over days is thin: QuantumNAS and the learned-mitigation study both
mention that noise drifts \cite{wang2022quantumnas,liao2024mlqem}, yet the multi-day hardware
study we found uses QAOA \cite{qaoa2026drift}, and the hardware evidence for conformal
prediction under time-varying noise comes from one day \cite{spencer2026aqcp}. Second,
corrections are mostly evaluated on one device at a time; the cross-device work we found
concerns noise models rather than classifier accuracy \cite{farib2026crossdevice}. Third, the
cost of a correction is usually described qualitatively, although on a free plan it decides what
can be run at all. Fourth, small hardware test sets make the accuracy of a single session
unreliable \cite{qbiobench2026}, so effects have to be pooled over a fixed test set and
repeated sessions.

To our knowledge, based on searches of OpenAlex, Semantic Scholar and arXiv in October and
November 2026, no earlier study repeats frozen quantum classifiers on more than one IBM Heron
processor over several days, separates drift from shot noise, compares several corrections by
their QPU cost, and checks conformal coverage on the same data. This is the gap that the
present work addresses.

% Drafted with an AI assistant on 12 Nov 2026 from the registered design (preregistration.md and
% config/experiment.toml). Revise it in your own words. Numbers are \val{} macros.
\section{Problem Statement}
\label{sec:problem}

\textbf{Problem.} A quantum classifier is trained in a noiseless simulator and then frozen. How
much do its accuracy and the reliability of its predictions change when the same circuits run
on real IBM Heron processors on different days and devices? How much of the change is genuine
drift rather than shot noise? And which low-cost corrections recover the loss for the least QPU
time? The gap in the literature that motivates these questions is summarised in
Section~\ref{sec:gap}.

\subsection{Setting}
Let $x\in[0,\pi]^{n}$ be an input with $n=\val{n-qubits}$ features and $\rho_\theta(x)$ the state
that a fixed circuit prepares on $n$ qubits. A model reads out Pauli expectation values and
passes them to a logistic head,
\begin{align}
\hat p(y=1\mid x)&=\sigma\!\Big(w^\top\frac{\phi(x)-\mu}{s}+b\Big), \label{eq:model}\\
\phi_j(x)&=\operatorname{tr}\!\big[P_j\,\rho_\theta(x)\big], \nonumber
\end{align}
where $P_j$ is a Pauli string and $\mu$, $s$ are the mean and standard deviation of the training
features. The circuit parameters $\theta$ and the head $(w,b,\mu,s)$ are fixed after training in
simulation. On device $d$ in session $t$, a job runs the circuits with $N=\val{shots}$ shots and
returns estimates $\hat\phi_{d,t}(x)$. Their expectation $\bar\phi_{d,t}(x)$ differs from
$\phi(x)$ because of gate, readout and decoherence errors, and it changes with $t$ as the
processor drifts and is recalibrated. Shot noise adds a variance of
$(1-\bar\phi_{d,t,j}^{2})/N$ to feature $j$. Accuracy is measured on a fixed test set of
\val{n-test} inputs, giving $A_{\mathrm{sim}}$ in the simulator and $A_{d,t}$ on hardware.

\subsection{Quantities to determine}
\begin{itemize}
\item \emph{Deployment gap and emulator error.} $G_{d,t}=A_{\mathrm{sim}}-A_{d,t}$ and
$E_{d,t}=A^{\mathrm{emu}}_{d,t}-A_{d,t}$, where $A^{\mathrm{emu}}_{d,t}$ is the accuracy
predicted by a noise model built from the calibration reported at submission.
\item \emph{Drift.} For each input and feature, the ratio of the variance of the estimate across
sessions to its expected shot variance,
\begin{equation}
r_d(x,j)=\frac{\operatorname{Var}_t\big[\hat\phi_{d,t,j}(x)\big]}
{\operatorname{E}_t\big[(1-\hat\phi_{d,t,j}(x)^{2})/N\big]}.
\label{eq:drift}
\end{equation}
A ratio of one means that shot noise alone explains the variation between sessions.
\item \emph{Noise map.} The scale $\gamma$ and shift $\beta$ in
$\bar\phi_{d,t,j}(x)\approx\gamma\,\phi_j(x)+\beta$, and their stability over $t$.
\item \emph{Repairs.} A repair $\mathcal R$ maps the hardware counts, and possibly $k$ labelled
hardware inputs, to corrected predictions. Its gain is $g(\mathcal R)=A^{\mathcal R}_{d,t}-
A_{d,t}$, and its cost $c(\mathcal R)$ is the extra QPU seconds per session.
\item \emph{Uncertainty.} Prediction sets $C_\alpha(x)\subseteq\{0,1\}$ with target coverage
$1-\alpha=\val{coverage-target-frac}$, judged by the empirical coverage
$\Pr[y\in C_\alpha(x)]$ and the mean set size.
\end{itemize}

\subsection{Research questions and hypotheses}
Table~\ref{tab:rq} links each research question to the quantity above and to the hypothesis that
was registered after the pilot runs and before the first campaign session (git tag
\texttt{prereg-v1}, \val{prereg-date}). Every hypothesis has a decision rule, fixed in advance and
based on 95\% bootstrap intervals, so that its outcome does not depend on the data.

\begin{table}[t]
\centering
\caption{Research questions, the quantities measured and the registered hypotheses.}
\label{tab:rq}
\footnotesize
\setlength{\tabcolsep}{3pt}
\begin{tabular}{@{}l>{\raggedright\arraybackslash}p{2.35cm}>{\raggedright\arraybackslash}p{4.2cm}@{}}
\toprule
 & Quantity & Registered hypothesis \\
\midrule
RQ1 & $G_{d,t}$, $E_{d,t}$ & H1: the emulator overestimates hardware accuracy ($E_{d,t}>0$) \\
RQ2 & drift ratio $r_d$; regression on chain calibration & H2: the ratio exceeds one. H4: readout error is the strongest predictor of feature error \\
RQ3 & $\gamma$, $\beta$ per session & none (descriptive) \\
RQ4 & $g(\mathcal R)$, $c(\mathcal R)$, coverage & H3: a 16-label refit beats normalisation. H5: split conformal calibrated in session 1 under-covers later sessions, while ACI holds \\
\bottomrule
\end{tabular}
\end{table}

\subsection{Constraints and scope}
The experiments use the free IBM Open Plan, which gives ten minutes of processor time per 28
days. The campaign is capped at \val{campaign-cap}~s of QPU time, jobs run when a shared queue
allows, and the models, datasets and qubit chains are fixed before the first session. We study
small classifiers on classical data (\val{n-qubits} qubits, two datasets) and do not train on
hardware. We do not claim quantum advantage: the quantum models are compared with tuned
classical baselines only to put their accuracy in context. The study is a measurement of
deployment behaviour on two processors of one family, not a general model of quantum noise.


\section*{Acknowledgment}
% Same wording as in main.tex; keep both documents consistent.
We acknowledge the use of IBM Quantum services for this work. The views expressed are those of
the authors, and do not reflect the official policy or position of IBM or the IBM Quantum team.
The analysis software and first drafts of parts of this manuscript were prepared with the help
of an AI assistant (Kiro); the author reviewed, revised and verified all code, text and results.

\bibliographystyle{IEEEtran}
\bibliography{generated/refs}

\end{document}
